\section{All weight-one Frame-shape forms}
The weight-one list separates a small cuspidal family from a larger
Eisenstein family, and requires the factors at primes dividing the
level to be displayed.
For a Frame shape \(\pi=\prod_d d^{r_d}\), define
\(\eta_\pi(\tau)=\prod_d\eta(d\tau)^{r_d}\). Its weight is
\(\tfrac12\sum_d r_d\).
Kond\=o's complete table of the Leech automorphism group gives
29 class entries of weight one, with 28 distinct Frame shapes
\cite[Table~I, pp.~355--357]{Kondo1985}.
The signs on class names below follow his convention on p.~354.

Let \(d_1,d_2\) be fundamental discriminants, allowing \(1\) for
the trivial character. Assume their characters have opposite parity:
\(\chi_{d_1}(-1)\chi_{d_2}(-1)=-1\). Write
\[
 E_{d_1,d_2}(\tau)=c_0+
 \sum_{n\ge1}\left(\sum_{a\mid n}
             \chi_{d_1}(a)\chi_{d_2}(n/a)\right)q^n.
\]
If neither character is trivial, \(c_0=0\); if one is trivial
and the other is \(\chi_d\), put \(c_0=L(0,\chi_d)/2\).
These are the weight-one Eisenstein series, with
\[
 L(E_{d_1,d_2},s)=L(s,\chi_{d_1})L(s,\chi_{d_2}),
\]
where the constant term is omitted in the Dirichlet transform.

\begin{theorem}\label{thm:weight-one-list}
The distinct weight-one eta products for \(M_{24}\) are
\[
 \eta(\tau)\eta(23\tau),\qquad
 \eta(3\tau)\eta(21\tau),\qquad
 \eta(12\tau)^2.
\]
They are irreducible dihedral cusp forms.
For \(\mathrm{Co}_0\), the five distinct cuspidal products are
those in Table~\ref{tab:weight-one-cusp}; all other 23 distinct
products are given in the following table.
For each noncuspidal row, if \(P(s)=\sum_t c_t t^{-s}\), then
\begin{equation}\label{eq:weight-one-eis}
 \eta_\pi(\tau)=\sum_t c_tE_{d_1,d_2}(t\tau),\qquad
 L(\eta_\pi,s)=P(s)L(s,\chi_{d_1})L(s,\chi_{d_2}).
\end{equation}
The primitive Galois representation in that row is
\(\chi_{d_1}\oplus\chi_{d_2}\), hence reducible.
Its Hecke character over the imaginary quadratic field of the
listed discriminant \(D\) is
\(\chi_{d_1}\circ\mathrm N_{K/\Q}\), equivalently
\(\chi_{d_2}\circ\mathrm N_{K/\Q}\).
The factors \(P\) are part of the displayed eta-product
\(L\)-series and cannot be discarded.
\end{theorem}

\begin{table}[htbp]
\centering
\caption{The five cuspidal weight-one products for \(\mathrm{Co}_0\).
The chosen CM realization is \(K=\Q(\sqrt{D_K})\), with
Hecke conductor \(f\mathcal O_K\).
The character is a faithful character of the indicated cyclic ring class
group, up to inversion.}
\label{tab:weight-one-cusp}
\begin{tabular}{@{}lrrrr@{}}
\toprule
Eta product&Level&\(D_K\)&\(f\)&Character order\\\midrule
\(\eta(\tau)\eta(23\tau)\)&23&\(-23\)&1&3\\
\(\eta(2\tau)\eta(22\tau)\)&44&\(-11\)&2&3\\
\(\eta(3\tau)\eta(21\tau)\)&63&\(-7\)&3&4\\
\(\eta(4\tau)\eta(20\tau)\)&80&\(-20\)&2&4\\
\(\eta(12\tau)^2\)&144&\(-4\)&6&4\\
\bottomrule
\end{tabular}
\end{table}

{\small
\begin{longtable}{@{}llrcll@{}}
\caption{The noncuspidal weight-one Frame shapes. Here \(X_p=p^{-s}\).}\\
\toprule Class&Frame shape&Level&\(\chi\)&\((d_1,d_2)\)&\(P(s)\)\\\midrule\endfirsthead
\toprule Class&Frame shape&Level&\(\chi\)&\((d_1,d_2)\)&\(P(s)\)\\\midrule\endhead
-6F & $\frac{1\cdot6^{6}}{2^{2}\cdot3^{3}}$ & 6 & $\chi_{-3}$ & $(1,-3)$ & $1 - X_{2}$ \\
8D & $\frac{8^{4}}{4^{2}}$ & 8 & $\chi_{-4}$ & $(1,-4)$ & $1 - X_{2}$ \\
+9B & $\frac{9^{3}}{3}$ & 9 & $\chi_{-3}$ & $(1,-3)$ & $1 - X_{3}$ \\
-12D & $\frac{2\cdot3^{3}\cdot12^{3}}{1\cdot4\cdot6^{3}}$ & 12 & $\chi_{-3}$ & $(1,-3)$ & $- (X_{2} - 1) (2 X_{2} + 1)$ \\
12G & $\frac{4^{2}\cdot12^{2}}{2\cdot6}$ & 12 & $\chi_{-3}$ & $(1,-3)$ & $- (X_{2} - 1) (X_{2} + 1)$ \\
+12H & $\frac{2^{3}\cdot6\cdot12^{2}}{1\cdot3\cdot4^{2}}$ & 12 & $\chi_{-4}$ & $(1,-4)$ & $1 - X_{3}$ \\
-12K & $\frac{1^{3}\cdot12^{3}}{2\cdot3\cdot4\cdot6}$ & 12 & $\chi_{-3}$ & $(1,-3)$ & $(X_{2} - 1) (2 X_{2} - 1)$ \\
-14B & $\frac{2^{2}\cdot14^{2}}{1\cdot7}$ & 14 & $\chi_{-7}$ & $(1,-7)$ & $1 - X_{2}$ \\
+15E & $\frac{1^{2}\cdot15^{2}}{3\cdot5}$ & 15 & $\chi_{-15}$ & $(-3,5)$ & $1$ \\
16A & $\frac{2^{2}\cdot16^{2}}{4\cdot8}$ & 32 & $\chi_{-8}$ & $(-4,8)$ & $1$ \\
-16B & $\frac{1^{2}\cdot16^{2}}{2\cdot8}$ & 16 & $\chi_{-4}$ & $(1,-4)$ & $(X_{2} - 1) (2 X_{2} - 1)$ \\
-18B & $\frac{1^{2}\cdot9\cdot18}{2\cdot3}$ & 18 & $\chi_{-3}$ & $(1,-3)$ & $(2 X_{2} - 1) (X_{3} - 1)$ \\
+18C & $\frac{1\cdot2\cdot18^{2}}{6\cdot9}$ & 18 & $\chi_{-3}$ & $(1,-3)$ & $(X_{2} - 1) (3 X_{3} - 1)$ \\
-18C & $\frac{2^{2}\cdot9\cdot18}{1\cdot6}$ & 18 & $\chi_{-3}$ & $(1,-3)$ & $- (X_{2} + 1) (X_{3} - 1)$ \\
+20C & $\frac{1\cdot2\cdot10\cdot20}{4\cdot5}$ & 20 & $\chi_{-20}$ & $(-4,5)$ & $1$ \\
-20C & $\frac{2^{2}\cdot5\cdot20}{1\cdot4}$ & 20 & $\chi_{-4}$ & $(1,-4)$ & $1 - X_{5}$ \\
24E & $\frac{2\cdot6\cdot8\cdot24}{4\cdot12}$ & 48 & $\chi_{-3}$ & $(-4,12)$ & $1$ \\
+24F & $\frac{1\cdot4\cdot6\cdot24}{3\cdot8}$ & 24 & $\chi_{-24}$ & $(-3,8)$ & $1$ \\
-24F & $\frac{2\cdot3\cdot4\cdot24}{1\cdot8}$ & 24 & $\chi_{-8}$ & $(1,-8)$ & $1 - X_{3}$ \\
-28A & $\frac{1\cdot4\cdot7\cdot28}{2\cdot14}$ & 28 & $\chi_{-7}$ & $(1,-7)$ & $(X_{2} - 1) (2 X_{2} - 1)$ \\
+30D & $\frac{1\cdot6\cdot10\cdot15}{3\cdot5}$ & 30 & $\chi_{-15}$ & $(-3,5)$ & $X_{2} + 1$ \\
-30D & $\frac{2\cdot3\cdot5\cdot30}{1\cdot15}$ & 30 & $\chi_{-15}$ & $(1,-15)$ & $1 - X_{2}$ \\
-30E & $\frac{2\cdot3\cdot5\cdot30}{6\cdot10}$ & 30 & $\chi_{-15}$ & $(-3,5)$ & $2 X_{2} + 1$ \\
\bottomrule\end{longtable}

}

\begin{proof}
For completeness, the finite classification and the identities can
be checked by short exact computations, with no identification from
approximate prime coefficients.
Selecting \(\sum r_d=2\) from the three pages of Kond\=o's
table gives exactly the rows displayed here. The only duplicate
Frame shape is \(1\cdot23\), for \(+23A,+23B\).
The \(M_{24}\) cycle-shape table selects the three products in
the statement \cite[Table~I, p.~154]{Koike1985}.

For every row, \(\sum d r_d=24\), so the eta quotient begins
with \(q\). With the level \(N\) in the table, its second eta
congruence is
\(\sum (N/d)r_d\equiv0\pmod{24}\).
Its character is the Kronecker character of the square class
\(-\prod d^{r_d}\), giving exactly the displayed \(\chi_D\).
For a cusp with denominator \(c\mid N\), the order is
\begin{equation}\label{eq:eta-cusp-order-weight1}
 \frac{N}{24\,c\,\gcd(c,N/c)}
       \sum_d\frac{\gcd(c,d)^2}{d}r_d.
\end{equation}
Substitution gives nonnegative orders for every row; all orders
are positive exactly for the five products in
Table~\ref{tab:weight-one-cusp}.
Thus the eta quotients and their proposed Eisenstein expressions
are holomorphic forms of the same weight, level and character.

We recall why a finite coefficient comparison proves these identities.
For a holomorphic form of weight \(k\) on \(\Gamma_0(N)\), the
valence formula applied to the product of its translates to
\(\mathrm{SL}_2(\Z)\) bounds its order at infinity by
\(k[\mathrm{SL}_2(\Z):\Gamma_0(N)]/12\), unless it is zero.
A power of that product removes its finite character if necessary.
Consequently two such forms agree if their coefficients through
\[
 B=\left\lfloor\frac{kN}{12}\prod_{p\mid N}(1+1/p)\right\rfloor
\]
agree. For the noncuspidal rows, the coefficients in
Table~\ref{tab:weight-one-checks} follow on the eta side by
multiplying \(\prod_{d,n}(1-q^{dn})^{r_d}\), and on the
Eisenstein side by the divisor sum in its definition.
The constant terms vanish: either both characters are nontrivial,
or \(\sum c_t=0\). This proves \eqref{eq:weight-one-eis}.
It also makes the character and level information directly
checkable.

Here is a separate exact identification of the five cusp forms.
For \(Q=[a,b,c]\), put
\(\theta_Q=\sum_{x,y\in\Z}q^{ax^2+bxy+cy^2}\).
The following identities hold:
\[
\begin{array}{c|cc|r}
 N&Q_0&Q_1&B\\\hline
23 &[1,1,6]&[2,1,3]&2\\
44 &[1,0,11]&[3,2,4]&6\\
63 &[1,1,16]&[4,1,4]&8\\
80 &[1,0,20]&[4,0,5]&12\\
144&[1,0,36]&[4,0,9]&24
\end{array}
\qquad
 \eta_\pi=\frac{\theta_{Q_0}-\theta_{Q_1}}2.
\]
Binary theta modularity follows from Poisson summation, with
level \(N\) and character \(\chi_{-N}\). The two sides have
the following common initial polynomials, with terms beyond
the indicated \(B\) omitted:
\[
\begin{array}{c|l}
23&q-q^2\\
44&q-q^3-q^5\\
63&q-q^4-q^7\\
80&q-q^5-q^9\\
144&q-2q^{13}
\end{array}
\]
These are obtained by enumerating the pairs \((x,y)\) in the
two positive quadratic forms; the same valence bound proves
all five identities.

Reduction of primitive positive binary quadratic forms gives
the complete class lists
\[
\begin{array}{c|l}
-23 &[1,1,6],\ [2,\pm1,3]\\
-44 &[1,0,11],\ [3,\pm2,4]\\
-63 &[1,1,16],\ [2,\pm1,8],\ [4,1,4]\\
-80 &[1,0,20],\ [3,\pm2,7],\ [4,0,5]\\
-144&[1,0,36],\ [5,\pm4,8],\ [4,0,9].
\end{array}
\]
The first two class groups have order three. In each of the
last three, only the principal class and the final class are
self-inverse, so the group of order four is cyclic.
For a faithful character \(\xi\) of one of these groups,
counting ideals of each class gives
\(\theta_\xi=\tfrac12\sum_{[Q]}\xi([Q])\theta_Q\).
Inverse classes have the same theta series. In order three their
two character values sum to \(-1\); in order four their values
cancel, and the order-two class has value \(-1\).
Thus \(\theta_\xi\) is precisely the displayed theta difference.
The characters are primitive at the conductors in the table:
the class groups of every proper divisor conductor have order
smaller than the stated character order. Since conjugation
sends \(\xi\) to \(\xi^{-1}\ne\xi\), their induced
two-dimensional representations are irreducible.
This proves the cuspidal and Hecke-character assertions.

For an Eisenstein row, let \(\chi_K=\chi_D\).
The table gives \(\chi_{d_2}=\chi_{d_1}\chi_K\) as primitive
characters away from their conductors. Restricting to \(G_K\)
identifies the two characters. Quadratic induction therefore
decomposes as
\[
 \operatorname{Ind}_{G_K}^{G_\Q}
        (\chi_{d_1}|_{G_K})
   =\chi_{d_1}\oplus\chi_{d_2}.
\]
The ideal-theoretic character is
\(\chi_{d_1}\circ\mathrm N_{K/\Q}\).
Its conductor norm is
\(|d_1d_2|/|D|\), which is \(1\) except for the row \(16A\)
(norm \(4\)) and \(24E\) (norm \(16\)).
The corresponding conductor ideals are \((2)\) in
\(\Q(\sqrt{-2})\), and \((4)\) in \(\Q(\sqrt{-3})\).
This proves the final assertion, including the ramified cases.
\end{proof}

{\small
\begingroup
\renewcommand{\tablename}{Table}
\begin{longtable}{@{}lrl@{}}
\caption{Exact coefficients through the weight-one Sturm bound \(B\).
The constant coefficient is zero in every row.}\label{tab:weight-one-checks}\\
\toprule Class&\(B\)&\((a_1,\ldots,a_B)\)\\\midrule\endfirsthead
\toprule Class&\(B\)&\((a_1,\ldots,a_B)\)\\\midrule\endhead
-6F & 1 & $(1)$ \\
8D & 1 & $(1)$ \\
+9B & 1 & $(1)$ \\
-12D & 2 & $(1,1)$ \\
12G & 2 & $(1,0)$ \\
+12H & 2 & $(1,1)$ \\
-12K & 2 & $(1,-3)$ \\
-14B & 2 & $(1,1)$ \\
+15E & 2 & $(1,-2)$ \\
16A & 4 & $(1,0,-2,0)$ \\
-16B & 2 & $(1,-2)$ \\
-18B & 3 & $(1,-2,0)$ \\
+18C & 3 & $(1,-1,-2)$ \\
-18C & 3 & $(1,1,0)$ \\
+20C & 3 & $(1,-1,-2)$ \\
-20C & 3 & $(1,1,0)$ \\
24E & 8 & $(1,0,-1,0,0,0,-2,0)$ \\
+24F & 4 & $(1,-1,-1,1)$ \\
-24F & 4 & $(1,1,1,1)$ \\
-28A & 4 & $(1,-1,0,-1)$ \\
+30D & 6 & $(1,-1,-1,1,-1,1)$ \\
-30D & 6 & $(1,1,1,1,1,1)$ \\
-30E & 6 & $(1,0,-1,-1,-1,0)$ \\
\bottomrule\end{longtable}

\endgroup
}

The qualification at bad primes is visible already in the row \(-12K\).
Its coefficients satisfy \(a(2)=-3\) and \(a(4)=3\), so it is
not an eigenfunction of \(U_2\): the required relation would be
\(a(4)=a(2)^2=9\). It is an Eisenstein eigenfunction at every
prime not dividing the level. Thus the assertion that every
\(\mathrm{Co}_0\) Frame-shape eta quotient is a simultaneous
newform is false.

For the reducible rows, Alp\"oge's theorem puts all nontrivial
zeros of the two primitive Dirichlet factors on \(\Re s=1/2\).
The factors \(P\) in the table also have zeros, on \(\Re s=0\)
or \(\Re s=1\), with possible cancellation at a pole of a
Dirichlet factor. For example, \(-12K\) has
\[
 P(s)=(1-2^{-s})(1-2^{1-s}),
\]
so it has infinitely many elementary zeros on \(\Re s=1\);
the zero at \(s=1\) itself cancels the zeta pole.
The five primitive dihedral cusp-form \(L\)-functions are
Hecke \(L\)-functions over quadratic fields, not products of
Dirichlet \(L\)-functions. Their Riemann hypotheses do not
follow from the Dirichlet theorem.
